\section{Why Stationary Laws Cannot Gather at an Unstable Optimum}\label{app:repulsion}

Use the bounded optimizer state and corrected recursion of
Appendix~\ref{app:transfer}. At zero parameters, each coordinate update
has mean zero. Curvature may depend on noise magnitudes; flipping any
individual noise-coordinate sign still leaves the joint law unchanged.

\begin{lemma}[Uniform conditional origin excitation]
\label{lem:uniform-origin-excitation}
Let $|e_i|\leq R_{\max}$, with $\E e_i^2>0$ for every coordinate.
For every old history and vector $v$, the origin conditional bracket obeys
\begin{equation}\label{eq:uniform-origin-bracket}
 v^\top q(0,Z)v\geq
 \left(\frac{1-\alpha}{R_{\max}+\epsilon}\right)^2
       \sum_i v_i^2\E e_i^2.
\end{equation}
Thus fresh noise remains present even after conditioning on the entire
past. Curvature may depend on noise; it does not enter the update at
zero parameters.
\end{lemma}
\begin{proof}
The averaged update at zero is zero. Conditional on future noise
magnitudes, every fresh sign has mean zero. For the update $k$ steps
ahead, averaging over fresh signs leaves only the old numerator. Thus
the $i$th corrector at zero has the exact representation
\[
 H_{0,i}(Z)=\sum_{l\geq0}e_{i,-l}a_{i,l}(|Z|),\qquad
 a_{i,l}(|Z|)=\sum_{k\geq0}w_{k+l}
       \E_{\rm future}\frac1{\sqrt{V_{i,k}}+\epsilon}.
\]
Here $|Z|$ denotes the coordinate-noise magnitude history. Its coefficients
are sign independent and satisfy
\[
 0\leq a_{i,l}\leq\frac{\alpha^l}{\epsilon},\qquad
 a_{i,0}\geq\frac{w_0}{R_{\max}+\epsilon}>0.
\]
The upper bound gives uniform absolute summability, at most
$R_{\max}/[\epsilon(1-\alpha)]$, justifying all exchanges. The lower
bound retains the $k=0$ term and uses $V_{i,k}\leq R_{\max}^2$.

Condition on the old history and the complete fresh magnitude vector
$r=(|e_{i,n+1}|)_i$. Then
\[
 H_{0,i}(Z_{n+1})=e_{i,n+1}a_{i,0}(|Z_{n+1}|)+b_i,
\]
where the coefficient and $b_i$ are measurable under this conditioning.
Because each coordinate sign can be flipped without changing the law,
the nonzero fresh signs are independent and equally likely positive or
negative conditional on $r$. Their cross covariances are zero, while
zero coordinates contribute zero. Therefore
\[
 \operatorname{Var}\!\left(v^\top H_0(Z_{n+1})\mid\mathcal F_n,r\right)
   =\sum_i v_i^2r_i^2a_{i,0}^2.
\]
The law of total variance, the coefficient lower bound and freshness of
$r$ prove \eqref{eq:uniform-origin-bracket}. Subtracting the past-measurable term
$\mathsf P H_0(Z_n)$ does not change this conditional variance.
\qed\end{proof}

\begin{proposition}[Stationary mass cannot concentrate at the repelling origin]
\label{prop:matrix-stationary-repulsion}
In dimension two assume the noise hypotheses of
Lemma~\ref{lem:uniform-origin-excitation}, and let $v=(1,-1)/\sqrt2$.
Suppose swapping the two coordinates swaps the averaged drift in the same
way (exchange equivariance), and its derivative at zero has a strictly negative
eigenvalue on $v$. Under the bounded native setup above, there exist
$\delta,c,\eta_0>0$ such that every invariant native-cone law with
$0<\eta<\eta_0$ satisfies
\[
 \Pp(\|u\|\geq\delta)\geq c,\qquad
 \E L_{\rm pop}(u)\geq
       \tfrac12\lambda_{\min}(\bar H)\delta^2c>0,
 \quad \bar H=\E A.
\]
The constants and complete sample law are fixed before $\eta$ decreases.
\end{proposition}
\begin{proof}
Write $y=v^\top x$ and $G_v=v^\top G$. Coordinate-swap symmetry implies
$G_v(z,0)=0$ in symmetric/antisymmetric coordinates. Integrating its
$y$ derivative gives $|G_v(x)|\leq L|y|$ on an enlarged bounded parameter box.
The strictly negative derivative at zero gives $a_0,\delta_0>0$ with
\[
 yG_v(x)\leq-a_0y^2\quad\hbox{when }\|x\|\leq\delta_0.
\]
The conditional noise bound in Lemma~\ref{lem:uniform-origin-excitation}
and the bracket's uniform Lipschitz bound imply that some fixed
neighborhood has
\[
 q_v(x,Z):=v^\top q(x,Z)v\geq q_{\min}/2>0.
\]
Globally on the enlarged box $0\leq q_v\leq q_{\max}<\infty$.
Set $h=6q_{\max}/a_0$ and
$\phi_\eta(y)=\log(y^2+h\eta)$. This test is bounded on the parameter
box for each positive $\eta$. Its first three derivative bounds are
$O(\eta^{-1/2})$, $O(\eta^{-1})$, $O(\eta^{-3/2})$.

Expand \eqref{eq:corrected-core} by Taylor's formula and use stationarity
to set the expected change of the test to zero. This gives
\begin{equation}\label{eq:conditional-log-generator}
 0=\E B_\eta(X,Z)+O(\sqrt\eta),\qquad
 B_\eta(x,Z)=-\phi_\eta'(y)G_v(x)
             +\tfrac\eta2\phi_\eta''(y)q_v(x,Z).
\end{equation}
Indeed, the first martingale term and its drift cross term have zero
expectation. The $R$ term times $\phi'$ and the third-order jump remainder
are $O(\eta^{3/2})$ before division by $\eta$. The drift-square term is
$O(\eta^2)$ because $|G_v(x)|\leq L|y|$ and
$|\phi_\eta''(y)|y^2\leq C$. Replacing $q(u,Z)$ by $q(X,Z)$ costs
$O(\eta^2)$ by $\|u-X\|\leq C\eta$. Throughout the calculation the variance remains conditional on the past;
it has not been replaced by its unconditional mean in
\eqref{eq:conditional-log-generator}.

On a fixed neighborhood where the drift and bracket bounds both hold,
put $z=y^2/(h\eta)$. Then
\[
 B_\eta(x,Z)\geq\frac{2a_0z}{1+z}
       +\frac{q_v(x,Z)}h\frac{1-z}{(1+z)^2}.
\]
For $z\leq1/2$, the second term is at least $q_{\min}/(9h)$ and the
first is nonnegative. For $z>1/2$, the first is at least $2a_0/3$ and
the second at least $-q_{\max}/h=-a_0/6$. Consequently
\[
 B_\eta(x,Z)\geq\kappa:=
       \min\{q_{\min}/(9h),a_0/2\}>0
\]
on that neighborhood for every history. Throughout the enlarged bounded
box it is at least
$-C_1$, with $C_1=2L+q_{\max}/h$. Equation
\eqref{eq:conditional-log-generator} therefore forces a fixed positive
probability outside the neighborhood for all sufficiently small $\eta$. Since
$\|u-X\|\leq C\eta$, halving its radius transfers the probability
bound to $u$. Positive $\bar H$ gives the stated risk lower bound.
\qed\end{proof}

